\section{Retained revision-65 abstract and leading statements}
\label{app:v66-retained-front}
This appendix retains the preceding abstract and leading statements.
Their theorem labels, hypotheses, conclusions and proof references are
unchanged. The new count-uniform contact and physical-trace results are
stated at the beginning of the present article.
\subsection*{The preceding abstract}

We study the joint displacement, collision count and flight time of
actual returns in a triangular periodic Lorentz gas. A visible spectral
branch is dominated at imaginary frequency by the positive pressure
of the original occupation record. At an exact resonance their
quadratic terms agree. We prove that this flat contact makes the
squared drift negligible compared with the damping. Consequently the
entire arithmetic transition kernel has one Gaussian shape, centered
at the physical mean and with the physical covariance, uniformly
through changes of arithmetic type. Its finite phases and damped
residues are retained.

The complete fixed-return law is therefore approximated in total
variation on its untruncated mixed record space by a damped arithmetic
factor times the physical return Gaussian. The same reference applies
to the graph-coupled collision and actual-return bridges, and to
record-only observations with their full selection events. The pressure
principle is stated in arbitrary dimension with separate physical
hypotheses and is also verified in a deterministic singular hyperbolic
baker map. We additionally verify the pressure hypotheses on a
correlated Markov baker map with a continuous deterministic roof. Its
exact stable coarea formula yields a pointwise roof local limit and
positive boundary-source height bounds uniformly throughout a nonzero
perturbation interval. Keeping two integer counts before scalar projection
resolves the arithmetic transition. For fixed endpoint witnesses the
critical density profile has exactly the damping predicted by pressure
contact. A separate positive coarea criterion states the exact-label
inputs needed for source heights. These independent realizations
do not replace the Lorentz record.

For the original Lorentz source, the endpoint contact determinant also
cancels a product of inverse incidences at every finite count. This yields
an incidence-source height gain. A two-coordinate physical coarea formula
bounds transverse clearance sources and retains their rank-deficient
part as a positive remainder. The count constants in these estimates
are exponential and are not identified with central-scale bounds.
Finite-type coarea further controls clearance tangent to a regular
roof fiber, including second-order folds. On an incidence-capped
physical graph, the remaining seam/rank caustic has finite fibers in
parameter--roof space. We localize the original positive clearance
source near that caustic, with its count-dependent constants retained.
A buffered algebraic caustic gives explicit separation and tube moduli
in both collision count and incidence cap. Fixed-block real quantifier
elimination supplies these finite-count bounds.

For normal-to-normal physical words we identify the intrinsic quadratic
coefficient with a reversible square-root determinant. Direct angular
coarea gives the density inside roof-critical clearance caustics,
including multiple seams and exact section decisions, with a
width-uniform finite-jet error.

The unrestricted pointwise raw-density target remains distinct from
these integrated conclusions. It still requires essential-height
control of the two positive physical boundary sources. No source
mass estimate is identified with that height control.
\subsection*{The preceding leading statements}

\paragraph{The original record.}
Let $J_{n,R}=(K_{n,R},N_{n,R},T_{n,R})$ record $n$ actual returns
of the triangular Lorentz collision map to $Y_R^*$, under
$\nu_R^*=c^{-1}\nu|_{Y_R^*}$, where $c=91/(10000\pi)$.
The density $p_{n,R}(k,m,u)$ is with respect to counting measure in
$(k,m)\in\Z^2\times\{n,n+1,\ldots\}$ and Lebesgue measure in $u>0$.
No collision or return is removed. Write
\[
 V_{n,R}(k,m,u)=
       \frac{(k_1,k_2,m-n/c,u-n\bar\tau_R/c)}{\sqrt n}.
\]
The physical return covariance is $D_R$; the collision occupation
covariance is $\Omega_R=cL_RD_RL_R^{\mathsf T}$, with
$\det L_R=c$. The finite peak data in
\eqref{eq:v41-peak-data} retain the exact occupation arithmetic.

\begin{leadtheorem}[Flat contact and the common Gaussian shape]
\label{thm:intro-v58-contact}
For every retained peak put $v_j=\mu_j-\mu_R$ and
$h_j=\|\Omega_R-\operatorname{Re}H_j\|$. Uniformly near zero damping,
\[
 |v_j|\le C\{\sqrt{\kappa_jh_j}+\kappa_j^{2/3}\},\qquad
 \sup_{0<\kappa_j\le\delta}\frac{|v_j|^2}{\kappa_j}\longrightarrow0
                 \quad(\delta\downarrow0).
\]
Thus $\sqrt m\,v_j\to0$ uniformly on each set $m\kappa_j\le L<\infty$.
For the complete transition kernel, set
\[
 \Phi_{m,R}(y)=c^{-1}\operatorname{Re}
             \sum_j e^{-ia_j(R)\cdot y}z_j(R)^mB_j(R).
\]
There are $a>0$ and $\theta_m\to0$ such that
\[
 \sup_{R,y}e^{a|y-m\mu_R|^2/m}
 \left|\mathcal L_{m,R}(y)-\Phi_{m,R}(y)
           g_{\Omega_R}\left(\frac{y-m\mu_R}{\sqrt m}\right)\right|
                                                      \le\theta_m.
\]
All damping, phases, residues and partition amplitudes are retained.
The covariance and center in the Gaussian are the physical ones,
also through changes of arithmetic type.
\end{leadtheorem}
\begin{proof}
Apply Theorem~\ref{thm:v58-flat-contact},
Corollary~\ref{cor:v58-lorentz-contact} and
Theorem~\ref{thm:v58-common-transition}.
\end{proof}

\begin{leadtheorem}[The complete arithmetic Gaussian return law]
\label{thm:intro-v58-record}
Set
\[
 \alpha_{m,R}(k,n,u)=c^{-1}\Phi_{m,R}(k_1,k_2,u,n),\qquad
 q^\circ_{n,R}(k,m,u)=n^{-2}\alpha_{m,R}(k,n,u)
                                        g_{D_R}(V_{n,R}(k,m,u)).
\]
Uniformly in $R\in[0.45,0.47]$, as the fixed return index $n\to\infty$,
\[
 \sum_{m=n}^\infty\sum_{k\in\Z^2}\int_0^\infty
                 |p_{n,R}(k,m,u)-q^\circ_{n,R}(k,m,u)|\,du\to0.
\]
The signed reference has vanishing negative mass and its positive
part has mass tending to one. Its normalized positive part
$\mathsf A^\circ_{n,R}$ approximates the original record law in
probability total variation. At fixed radius the damped factor may
be replaced by the exact finite residue $\mathfrak a_R(k,n,m)$,
including zero classes.

For the actual collision and return bridges from the same trajectory,
\[
 \left\|\operatorname{Law}_{\nu_R^*}
 (J_{n,R},\mathcal B_{N_{n,R},R},\mathcal Y_{n,R})
      -\mathsf A^\circ_{n,R}\otimes
       \operatorname{Law}(\mathbb B_{\Omega_R},
                    c^{-1/2}L_R^{-1}\mathbb B_{\Omega_R})
                       \right\|_{\mathrm{obs},\mathrm{BL}}\to0.
\]
The norm is variation in the observation and bounded-Lipschitz dual
in the two paths, not path-space total variation. The record becomes
independent of this coupled bridge pair; its two paths are not
independently sampled. A record-only observation kernel contracts the
error. If that error is $\Delta_n^\circ$, selection on its complete
output event of reference mass at least $r_n>\Delta_n^\circ$ costs
at most $\Delta_n^\circ/(r_n-\Delta_n^\circ)$ in probability TV and
twice that amount in the joint path-dual norm.
\end{leadtheorem}
\begin{proof}
Theorem~\ref{thm:v58-canonical-raw} and
Corollary~\ref{cor:v58-canonical-bridge} prove these statements.
Probability TV is one half of variation mass.
\end{proof}

\paragraph{Why the centers can be removed.}
The positive pressure dominates a visible centered branch at an
imaginary tilt. Their difference is nonnegative, has value $\kappa_j$
and gradient $v_j$, and its Hessian at zero is
$\Omega_R-\operatorname{Re}H_j$. At zero damping this Hessian
vanishes. Retaining it in Taylor's formula gives the first theorem;
a bound by an unspecified quadratic constant gives only
$|v_j|^2\le C\kappa_j$. The new proof also gives an explicit modulus
in the covariance mismatch, but does not assert a polynomial rate
in the collision count. It uses no derivative in the table parameter
and no spectral estimate at a growing Fourier band.

\paragraph{A reusable statement and a second realization.}
Section~\ref{sec:v58-pressure-contact} separates positive physical
moments, branch visibility and second-order contact as hypotheses of
a scalar theorem in arbitrary dimension. They are verified for the
original occupation spectrum and, independently, for a deterministic
three-strip baker map. In the latter example a peak has
$\kappa_\epsilon=\pi^2\epsilon^2/9+O(\epsilon^3)$ but
$v_\epsilon=\pi^2\epsilon^3/18+O(\epsilon^4)$. This is a second
singular hyperbolic realization of the spectral principle, not a
substitute for the Lorentz density theorem. The baker observable has
finite support; no continuous density theorem is asserted for it.

\paragraph{A correlated roof and a direct height calculation.}
Sections~\ref{sec:v59-markov-realization} and~\ref{sec:v59-stable-coarea}
verify the pressure/contact inputs in a singular Markov baker map whose
state correlations do not vanish. Its positive roof is a finite-state
reward plus the actual stable endpoint difference, not an added random
smoothing. For its unperturbed roof density $p_m^{\rm M}$, with
$\bar\tau=24/7$ and $\sigma^2=204/343$, we prove
\[
 \operatorname*{ess\,sup}_{t\in\R}
 \left|\sqrt m\,p_m^{\rm M}(t)
       -g_{\sigma^2}((t-m\bar\tau)/\sqrt m)\right|
                 \le C\log(2+m)/\sqrt m.
\]
The three stated boundary sources have normalized heights bounded by
$C\{s+(3/4)^{\lfloor m/2\rfloor}\}$ at zero perturbation. Those results
remain in Section~\ref{sec:v59-stable-coarea}. The following theorem
places density inversion, source heights and the moving peak on the
same perturbed family. It uses the original two integer counts before
projection, not a presumed common scalar lattice span.

\begin{leadtheorem}[Pointwise inversion through a perturbed arithmetic family]
\label{thm:intro-v60-perturbed}
On the Markov baker with transition matrix
$P=\left(\begin{smallmatrix}3/4&1/4\\1/3&2/3\end{smallmatrix}\right)$
and invariant state weights $\pi=(4/7,3/7)$, retain the positive roof
$\tau_\epsilon=3+I_1+\epsilon I_0I_1+y_1-y_0$.
For every fixed $0<\epsilon_0<1/4$, uniformly for
$|\epsilon|\le\epsilon_0$,
\[
 \operatorname*{ess\,sup}_{t\in\R}
 \left|\sqrt m\,p_{m,\epsilon}(t)
 -g_{\sigma_\epsilon^2}((t-m\bar\tau_\epsilon)/\sqrt m)\right|
 \le C\frac{\log(2+m)^{3/2}}{\sqrt m},
\]
where $\bar\tau_\epsilon=(24+2\epsilon)/7$ and
$\sigma_\epsilon^2=(204+384\epsilon+198\epsilon^2)/343$.
For the three original boundary sources,
\[
 \sqrt m\left(\|p_{m,\epsilon,s}^{\rm vert}\|_\infty
 +\|p_{m,\epsilon,s}^{\rm stab}\|_\infty
 +\|p_{m,\epsilon,s}^{\rm age}\|_\infty\right)
 \le C\{s+(3/4)^{L_m}\},\quad
 L_m=\min(\lfloor m/2\rfloor,\lfloor\sqrt m\rfloor).
\]
Lipschitz endpoint weights retain the explicit arithmetic profile
\eqref{eq:v60-arithmetic-profile}. On the critical scale
$\epsilon=c_m/\sqrt m$, bounded $c_m$, the fixed witness pair
$e^{2\pi iy_0},e^{-2\pi iy_m}$ has leading density
\[
 g_{204/343}(z)
 e^{-2\pi i(r_\epsilon-(16/17)c_mz)}e^{-12\pi^2c_m^2/119},
 \quad z=\frac{t-m\bar\tau_\epsilon}{\sqrt m},\quad
 r_\epsilon=t-3m-2m\epsilon/7,
\]
with the same normalized error order. In particular the damping
coefficient agrees with the moving pressure peak on this actual source.
\end{leadtheorem}
\begin{proof}
Apply Theorems~\ref{thm:v60-perturbed-density},
\ref{thm:v60-perturbed-height} and~\ref{thm:v60-critical-profile}.
\end{proof}

Finite-state Markov additive density theory and compact-family versions
already exist \cite{HL,HLuniform}. The new conclusion in this proof
chain concerns exact coarea through the scalar arithmetic transition,
positive source heights and the visible critical endpoint profile.
It does not assert a Lorentz height bound by analogy.

\paragraph{Density estimates on the original physical boundary source.}
For a regular $m$-collision orbit set
$\mathcal I_{m,R}=\prod_{j=0}^m\cos(\varphi\circ T_R^j)^{-1}$.
Theorem~\ref{thm:v61-inverse-incidence-density} proves that this unbounded
insertion has a roof density at most $CA^m$, for fixed $C>0$, $A>1$.
The same determinant controls all contact incidences simultaneously.
For the original first-incidence source, and for the part of the original
clearance source on which
$|\det D_{(\alpha,p)}(L_{m,R},Z_d\circ T_R^{j+1})|\ge\eta$,
Corollary~\ref{cor:v61-incidence-height-gain} and
Theorem~\ref{thm:v61-transverse-clearance} give
\[
 \operatorname*{ess\,sup}_u\sum_{n=1}^m\sum_k
 b^{\varepsilon,\mathrm{inc},1}_{n,k,m,R}(u)\le CA^m\varepsilon,
 \qquad
 \operatorname*{ess\,sup}_u\sum_{n=1}^m\sum_k
 b^{\varepsilon,\mathrm{tr},\eta,1}_{n,k,m,R}(u)
       \le CA^m\varepsilon/\eta.
\]
These estimates allow count-dependent thresholds and bounded source
insertions. They are derived on the true circular collision geometry,
not on a Markov replacement. They do not bound the rank-deficient
clearance source or remove the exponential count factor.

\paragraph{The unrestricted raw target and the proof route.}
The paper retains its four-coordinate record and unrestricted
two-sided pointwise target. The positive incidence and clearance
sources in Corollary~\ref{cor:v51-positive-criterion} still require
essential-height bounds at the exact labels. Removing the reference's
moving centers does not remove those physical sources. The new
reference positivity statement likewise does not give a pointwise
upper bound for the true density. No unrestricted same-roof bridge
or forward essential likelihood conclusion is drawn from the
integrated theorem.

The direct physical source calculation is in
Sections~\ref{sec:v61-inverse-incidence} and~\ref{sec:v61-clearance-coarea}.
Its inputs are the endpoint action and collision-graph complexity of
Section~\ref{sec:v44-complete-height}, the exact contact determinant of
Section~\ref{sec:v45-critical-collars}, and the next-collision clearance
geometry of Section~\ref{sec:v49-physical-layers}. Thus its short proof
route does not require the independent Markov realization.
The positive exact-label criterion is in
Section~\ref{sec:v60-positive-coarea}. Its perturbed verification and
critical density profile are in Sections~\ref{sec:v60-perturbed-density}
and~\ref{sec:v60-heights-contact}; the required two-count Fourier and
positive cylinder arguments are proved there. The scalar root,
finite-witness and complex Gaussian details are given
first in Section~\ref{sec:v59-contact-details}. The canonical Lorentz
route then runs through Sections~\ref{sec:v58-pressure-contact}
and~\ref{sec:v58-common-raw-law}, followed by the full-record tail
and joint clock arguments of
Sections~\ref{sec:v57-transition-tails} and~\ref{sec:v57-coupled-bridges}.
They state their inherited inputs locally. The entire earlier proof
chain and all earlier statements remain compiled; the previous front
matter is in Appendix~\ref{app:v58-retained-front}.
The comparison with billiard endpoint and suspension local limits
\cite{DPZ,DN} concerns the precise observable and topology, not a
claim that pressure domination or Gaussian differentiation is new.
The additional conclusion here is the uniform common Gaussian shape
of a visible moving resonance family, proved before summing the
original exact-return law.

\paragraph{Finite-type clearance and its physical caustic.}
The first-clearance source is further resolved on the actual roof
fibers in Sections~\ref{sec:v62-finite-type-clearance} and
\ref{sec:v62-clearance-caustics}. If a coordinate derivative of the
roof is at least $a$ and the $r$th clearance derivative along its level
curve is at least $b$, its exact-label essential height is at most
$C_rA_r^m a^{-1}(\varepsilon/b)^{1/r}$. In particular the second-order
estimate includes fold contacts with vanishing two-coordinate
Jacobian. The complete complementary source remains positive.

For fixed $m$ and incidence cap $\chi>0$, the physical clearance
caustic in $(R,t)$ has at most $CA^m$ points in each fixed-radius
fiber. A compact semialgebraic separation inequality bounds the
original source away from its neighborhood and gives an exact
raw-error formula retaining only the caustic source and explicit
finite-count errors. The neighborhood also has a power measure bound,
which combines with the inherited weak-integrability estimate.
These are geometric reductions for the Lorentz source, not a transfer
of the Markov model theorem. Their count-dependent constants do not
establish the two ordered central-height limits.


\paragraph{Effective moduli on the physical graph.}
Sections~\ref{sec:v63-effective-caustics} and
\ref{sec:v63-positive-localization} quantify a buffered version of
this localization. For one fixed $C_0$, put
\[
 E_m=\left\lceil2^{C_0(m+1)^{16}}\right\rceil,\qquad K_m=2^{E_m}.
\]
The source cap is $\chi$, while the critical-value reference uses
$\chi/2$. Theorems~\ref{thm:v63-buffered-separation} and
\ref{thm:v63-effective-tube} give
\[
 (\chi d)^{E_m}\le K_m\bigl(\bigl||Z|-R\bigr|+|\mathcal J|\bigr),
 \qquad \sup_R|\text{vertical }h\text{-tube}|
                  \le K_m\chi^{-1}h^{1/E_m}.
\]
The reference is an explicitly enlarged weak-first-hit critical graph,
not a replacement source. The original first-clearance source is
restricted positively to its neighborhood; its complement has the
height bound in Theorem~\ref{thm:v63-positive-localization}. The cap
buffer makes the compactness argument valid while the cap varies.
Standard effective quantifier elimination supplies the algebraic
bounds, with its integer coefficient control. The new application is
the fixed-format physical graph and the exact positive source formula,
not a new quantifier-elimination theorem. The very large count budgets
and the remaining tube height are kept distinct from the ordered
pointwise endpoint.


\paragraph{Intrinsic weights and density inside the caustic.}
The physical caustic is not used solely as an exceptional roof set.
Sections~\ref{sec:v65-reversible-critical}--\ref{sec:v65-affine-caustic}
calculate its roof-critical contribution locally on the same collision
source. They retain normal endpoints, positive selected incidences,
all competing-hit restrictions and the original section decisions.
A boundary branch linearization is distinguished from the derivative
of the billiard map at a singular state.

\begin{leadtheorem}[Reversible weights and physical caustic profiles]
\label{thm:intro-v65-caustic}
For a normal-to-normal $m$-flight word, let $M_z$ be its canonical
half-word derivative, $S=\operatorname{diag}(1,-1)$ and
$P_z=S M_z^{-1}S M_z$. The full local quadratic coefficient for
$\nu/c_*$ is
\[
 J_z=\frac{1}{R c_*\sqrt{|\det(I-P_z)|}}
     =\frac{1}{R c_*}\frac{\lambda_z^{-1/2}}{1-\lambda_z^{-1}},
          \qquad\lambda_z>1.
\]
At a regular word $P_z=D T_R^{2m}(z)$; at a competing-hit boundary
it denotes the physical-side branch expression only. The original
labels remain those of the first $m$ flights.

At a simple clearance seam through a roof minimum $t_z$, the
original local strip $0<g<s$ has density
\[
 b_{z,s}(t_z+h)=\frac{J_z}{\pi}
     \arcsin\!\min\left\{1,\frac{s}{\kappa_z\sqrt{2h}}\right\}
          +O_z(J_zh^{1/4}),\qquad h>0,
\]
uniformly in the width-to-roof ratio. Here
$\kappa_z^2=dg(z)(\nabla^2F(z))^{-1}dg(z)^{\mathsf T}>0$.
In particular its physical one-sided coefficient is $J_z/2$.
For any fixed word with positive selected incidences, multiple
clearance seams and section-boundary coincidences are allowed in
the affine angular profile of Section~\ref{sec:v65-affine-caustic}.
For all sufficiently small widths relative only to the local
incidence cap and the fixed near-clearance scale, it satisfies
\[
 \sum_{n,k}\left|b_{z,n,k}^{\varepsilon,\mathrm{clr}}(t_z+h)
       -\frac{J_z}{2\pi}\mathcal A_{z,n,k}(\varepsilon,h)\right|
                    \le C_zJ_zh^{1/4},\qquad
          0\le\sum_{n,k}\mathcal A_{z,n,k}\le2\pi.
\]
All constants and neighborhoods in these profile assertions are local
finite-word data. No ordered long-count height conclusion is included
in their statements.
\end{leadtheorem}
\begin{proof}
Theorem~\ref{thm:v65-reversible-weight} and
Proposition~\ref{prop:v65-boundary-monodromy} give the coefficient.
Theorems~\ref{thm:v65-seam-profile} and
\ref{thm:v65-affine-label-profile} give the two local profiles.
\end{proof}

The determinant relation is an explicit two-dimensional application
of the action--monodromy principle, not a replacement for the
classical Hill theory \cite{BT65}. The density calculation supplies
the original positive source fraction and its exact-label sum.
Its affine offsets avoid presuming that the fixed reconstruction
width is smaller than every nonzero clearance margin. The ordered
trace transfer in Corollary~\ref{cor:v65-ordered-trace-transfer}
states the further concentration input explicitly. Neither that input
nor uniform coverage of all long words is supplied by a finite-count
profile, and the full pointwise arithmetic target remains unchanged.

